% =====================================================================
%  Capítulo 1 — O problema
% =====================================================================
\chapter{O problema}
\label{cap:problema}
\rubricaJustificativa

\begin{orientacao}
Este capítulo tem uma única função: convencer quem lê de que o sistema
descrito nos capítulos seguintes resolve um problema que existe fora da
sala de aula. Duas a três páginas bastam. Se a equipe não consegue
escrever este capítulo sem falar de sensor, placa ou protocolo, o
problema ainda não foi definido --- foi escolhida uma solução e
procurado um problema para ela.
\end{orientacao}

\section{Contexto}

\begin{orientacao}
Onde o problema acontece, com quem, e há quanto tempo. Se algum
integrante tem acesso real a esse contexto (a propriedade rural, o
laboratório, o galpão, a residência), diga isso aqui e diga com quem a
equipe conversou. Uma conversa de quinze minutos com quem convive com o
problema vale mais que três parágrafos de citação genérica sobre o
crescimento da IoT.
\end{orientacao}

Ambientes que dependem de faixas estreitas de temperatura e umidade ---
salas de manipulação farmacêutica, câmaras frias de alimentos, salas de
equipamentos de TI --- ainda são, em grande parte, monitorados à mão. A
equipe conversou com um associado que atua na área comercial e de
gestão logística de uma farmácia de manipulação, e esse estabelecimento
foi adotado como cenário de teste do projeto. Ali, o controle de
temperatura e umidade é exigido pela RDC 67/ANVISA e é hoje inteiramente
manual: os farmacêuticos responsáveis técnicos consultam termômetros
espalhados pela sala --- inclusive dentro do freezer, em alguns pontos
--- e anotam o horário e a leitura em uma folha fixada na parede. A
verificação ocorre a cada hora, em turnos de 12 horas, o que resulta em
doze checagens manuais por dia. Os registros são depois consolidados em
relatórios e arquivados em gavetas, para eventual auditoria.

O mesmo problema de fundo --- ausência de registro contínuo e de
sinalização automática de desvio --- se repete em outros ambientes
controlados, como câmaras frias de estabelecimentos alimentícios e
salas de equipamentos de TI, o que motivou a equipe a projetar uma
solução que não fosse específica de um único segmento.

\section{Pergunta de projeto}

\begin{orientacao}
A pergunta no formato \enquote{Como podemos\ldots?} que a equipe entregou
em 24/08, já refinada pela mentoria de 14/09. Uma frase, em destaque.
Ela precisa ser específica o bastante para que se possa dizer, ao fim do
ciclo, se foi respondida ou não.
\end{orientacao}

\begin{quote}
\textbf{Como podemos} capturar as informações de temperatura e umidade
de uma sala e enviar para um servidor central, que reúne os dados dos
nós distribuídos e mostra a situação atual do local, \textbf{para} os
responsáveis por ambientes críticos \textbf{de modo que} a condição de
cada sala seja conhecida e sinalizada sem depender de rondas manuais
\textbf{?}
\end{quote}

A pergunta evoluiu desde a proposta inicial, que falava em avisar os
responsáveis \enquote{pelo celular}. Ao considerar os aspectos de
infraestrutura, capacidade e limitação dos sensores disponíveis, a
equipe percebeu que o envio remoto pressupõe um nó de borda que já
funcione de forma confiável isoladamente --- e foi isso que o primeiro
ciclo (AP1) construiu. Este segundo ciclo responde à outra metade da
pergunta: o envio das leituras a um ponto central por uma
infraestrutura de comunicação, e a formação, a partir delas, do
conjunto de dados do projeto.

\section{Quem usa o sistema, e o que muda para essa pessoa}

\begin{orientacao}
Descreva a pessoa concreta que vai conviver com o protótipo. O que ela
faz hoje, quanto tempo gasta nisso, o que ela deixa de fazer quando o
sistema funcionar. Esta seção é a que transforma o projeto em requisito
--- a maior parte dos requisitos do Capítulo~\ref{cap:requisitos} deveria
ser derivável daqui.
\end{orientacao}

A usuária de referência é a farmacêutica responsável técnica pela sala
de manipulação, que responde legalmente pela conformidade do ambiente
perante a fiscalização da ANVISA e cuida hoje de uma farmácia mais uma
filial. Ela é metódica e desconfia de processo manual, mas é justamente
o processo manual --- termômetro de gancho, planilha, caderno de
controle e o alarme do próprio celular --- que ela tem à disposição.
Sua necessidade central não é visualizar um gráfico bonito: é poder
garantir, a qualquer momento e sem correria na véspera de uma
fiscalização, que a fórmula manipulada seguiu a faixa de temperatura
exigida. Hoje isso depende de ela lembrar de checar o termômetro doze
vezes por dia; com o sistema instalado, a sinalização de desvio deixa
de depender da sua ronda e passa a ser automática e imediata, no
próprio ambiente.

\section{Onde o nó fica, e onde está a internet}
\label{sec:local}

\begin{orientacao}
Novo na AP2. Diga onde o nó fica instalado e onde está o primeiro ponto
com internet (roteador, \emph{gateway}, torre de celular): é dessa
distância, e do que há entre os dois pontos (paredes, porta de câmara
fria, laje), que sai o requisito de alcance do
Capítulo~\ref{cap:requisitos} e a primeira linha da matriz de decisão
da Seção~\ref{sec:tecnologia}. Se a campanha de coleta foi feita em
outro lugar que não o cenário de referência, descreva os dois e diga a
diferença.
\end{orientacao}

No cenário de referência --- a farmácia de manipulação descrita na
Cap.~\ref{cap:problema} ---, os pontos que hoje recebem termômetro são
a própria sala de manipulação e o interior do freezer. O nó do
TermoLOG fica na parede da sala de manipulação, no lugar do termômetro
de gancho que a farmacêutica consulta hoje, longe da saída do
ar-condicionado; o freezer é o segundo ponto, e o mais difícil para o
rádio. O primeiro ponto com internet é o roteador WiFi do
estabelecimento. \textbf{A planta da farmácia não foi levantada}: a
posição do roteador, a distância até a sala e as paredes no caminho não
são conhecidas, pelo mesmo motivo que tirou a campanha de lá (acesso
restrito às áreas monitoradas). Por isso o requisito de alcance não sai
de uma distância medida no local, e sim de uma \textbf{regra de
instalação} tirada das medições da equipe: o nó pode ficar a até
\SI{8}{\metre} e uma parede de alvenaria do ponto de acesso, e o RSSI é
conferido no ponto antes de fixá-lo (\RNF{06}).

O que a equipe já mediu é quanto cada obstáculo custa em
\SI{2,4}{\giga\hertz}, com o próprio rádio do ESP32, em um ambiente
interno de teste (bateria de 24/10, Seção~\ref{sec:tecnologia}): duas
portas de madeira custaram cerca de \SI{2,6}{\deci\bel}, uma parede de
alvenaria cerca de \SI{6}{\deci\bel}, e o enlace deixou de ser
confiável a partir de duas paredes ou de \SI{11}{\metre}. Desses números
sai a regra de instalação do \RNF{06}: RSSI p10 de pelo menos
\SI{-86}{dBm} no ponto do nó, o p10 do último ponto sem perda
(\SI{8}{\metre} e uma parede), medido antes de fixá-lo. A atenuação da
porta e das paredes metálicas de um freezer ou câmara fria \textbf{não
foi medida}, e é o caso mais desfavorável do cenário de referência. É
por causa dele que o produto não usa \SI{2,4}{\giga\hertz} entre os nós:
os nós de borda falam LoRa com um nó central, que fica na tomada perto do
roteador (Cap.~\ref{cap:arquitetura}), e só o nó central precisa
alcançar a internet.

A campanha de coleta deste ciclo (Seção~\ref{sec:campanha}) \textbf{não
foi feita no cenário de referência}. O nó ficou em cima de uma
escrivaninha, num cômodo com ar-condicionado da residência de um
integrante (João Pedro Saleh), fora do fluxo de ar do aparelho, com o
roteador no mesmo ambiente: o RSSI ficou entre \SI{-15}{} e
\SI{-29}{dBm}, mais de \SI{55}{\deci\bel} acima da regra do \RNF{06}. O
que muda em relação à farmácia: não há paredes nem a porta metálica do
freezer no caminho do rádio, e o ambiente é um cômodo comum, cuja
temperatura segue o ar-condicionado e o calor da tarde, e não uma sala
com faixa controlada. A conectividade medida na campanha é, portanto, o
melhor caso de rádio; e o fenômeno medido, o desligamento do
ar-condicionado, é o análogo doméstico de uma falha de climatização.
A farmácia ficou de fora porque o acesso às áreas monitoradas é
restrito: testar lá exigiria sanitização do equipamento e permissão dos
proprietários, que não havia como obter dentro do prazo do ciclo.

\section{Escopo deste ciclo}

\begin{orientacao}
Duas listas curtas e honestas. A segunda é a que o docente lê com mais
atenção: equipe que não sabe dizer o que \emph{não} vai fazer costuma
não entregar nem o que disse que faria.
\end{orientacao}

\noindent\textbf{Está neste ciclo (\moduloDoCiclo):}
\begin{itemize}
  \item integração do nó da AP1 (ESP32~+~DHT22, filtragem e alarme local) a
    uma infraestrutura de comunicação, com a tecnologia escolhida e
    justificada na Seção~\ref{sec:tecnologia};
  \item publicação de cada leitura em um \emph{broker} MQTT, segundo um
    contrato de dados e de tópicos documentado
    (Seções~\ref{sec:contrato-dados} e~\ref{sec:topicos});
  \item registro persistente das leituras por um gravador, que forma o
    conjunto de dados do projeto --- o registro que a AP1 adiou para
    nascer já associado ao dado que trafega pela rede;
  \item uma campanha de coleta com o nó operando sem ninguém olhando,
    e a medição da qualidade do conjunto formado
    (Seções~\ref{sec:campanha} a~\ref{sec:dataset});
  \item a medição do custo real do rádio, que substitui a estimativa do
    orçamento de energia da AP1 (Seção~\ref{sec:custo-radio});
  \item a avaliação, em bancada, dos dois rádios do próprio ESP32 (ESP-NOW
    e BLE) nas mesmas distâncias e obstáculos, como alternativa ao LoRa
    enquanto o módulo comprado pela equipe não chegava; o resultado
    sustenta a escolha da Seção~\ref{sec:tecnologia};
\end{itemize}

\noindent As pendências herdadas da AP1 --- correção da alteração de
faixa (\RF{02}), redução da latência do alarme (\RNF{02}),
\emph{timeout} na leitura do sensor e log formal do \CT{04} --- estavam
previstas para este ciclo, mas \textbf{não foram tratadas}: o trabalho
do Módulo~2 foi para a comunicação, num \emph{firmware} separado do da
AP1, e elas continuam abertas (Cap.~\ref{cap:conclusao}).

\noindent\textbf{Não está neste ciclo, e por quê:}
\begin{itemize}
  \item \textbf{Treinamento e embarque de um modelo} --- é o objeto do
    Módulo~3; este ciclo entrega o conjunto de dados que o alimenta;
  \item \textbf{Notificação remota (aplicativo, push, WhatsApp) e
    acionamento automático de climatização} --- continuam na visão de
    produto de longo prazo; com a rede, passam a ser tecnicamente
    possíveis, mas não cabem no prazo deste ciclo;
  \item \textbf{Integração do nó ao painel central web e relatórios em
    PDF/Excel para auditoria} --- o servidor do produto (API em FastAPI,
    banco PostgreSQL com registros imutáveis e painel web com um card por
    sala) foi construído em paralelo pela equipe. A placa ainda não fala
    com ele, porque o nó central não existe em hardware
    (Cap.~\ref{cap:arquitetura}): neste ciclo, o dado da placa termina no
    gravador (Seção~\ref{sec:caminho}), e as leituras da campanha foram
    levadas ao servidor por uma ponte que faz o papel do nó central
    (Seção~\ref{sec:ponte}, \CT{16}). Relatórios em PDF/Excel continuam
    fora do ciclo;
  \item \textbf{Rede com vários nós} --- a campanha usa
    um nó (\texttt{eq01-no02}); a escala para um nó por sala, com enlace
    LoRa até um nó central, depende do módulo LoRa e fica para a AP3 em
    diante.
\end{itemize}

\section{Soluções existentes}

\begin{orientacao}
Duas ou três soluções que já resolvem algo parecido --- produtos
comerciais, projetos abertos ou trabalhos publicados --- e, para cada
uma, uma frase dizendo por que ela não serve ao contexto da
Cap.~\ref{cap:problema}. Preço, dependência de infraestrutura que não
existe no local e falta de acesso ao código são razões legítimas.
\enquote{Não encontramos nada parecido} quase nunca é verdade e
enfraquece o relatório. Cite pela norma: \citeonline{alfuqaha2015} é um
ponto de partida razoável para o panorama de arquiteturas e protocolos.
\end{orientacao}

O mercado de registradores de temperatura para ambientes controlados é
maduro, mas caro para o caso de uma farmácia de manipulação de pequeno
porte (preços levantados pela equipe na AP1, em setembro de 2026). O Novus LogBox Wi-Fi, datalogger industrial com nuvem própria,
custa cerca de R\$~3.266,00 por unidade; a sonda Wi-Fi Testo Saveris
2-H1 custa por volta de R\$~3.608,00 avulsa; e a plataforma em nuvem
SyOS, voltada à farmácia, exige contrato mínimo de R\$~18.000,00 (cinco
sensores por 24 meses). O concorrente de preço mais próximo, o
datalogger portátil Elitech RC-5+ (cerca de R\$~220,00), resolve o
registro, mas é offline: a leitura só é conhecida quando alguém conecta
o aparelho por USB, o que reproduz o mesmo atraso do processo manual
que a equipe já observa na farmácia. Nenhuma das soluções pesquisadas
--- incluindo a Sensorweb, que opera em comodato sem preço público ---
aciona uma ação corretiva por conta própria: todas se limitam a
monitorar e alertar. O TermoLOG, com um nó por sala de R\$~78,90 em
ESP32 e DHT22 (R\$~105,20 com todo o protótipo, Tabela~\ref{tab:bom}),
não compete em recursos de nuvem com essas
soluções neste ciclo, mas parte de um custo por ponto de medição cerca
de trinta vezes menor (cerca de vinte vezes, com o módulo LoRa do
produto), o que muda o cálculo de viabilidade para
farmácias de pequeno e médio porte.
